\chapter{Aeronavegabilidad inicial y continuada}\label{ch:16}
\objetivos{Separar certificado de aeronavegabilidad, revisión ARC y certificación de mantenimiento; interpretar la revisión conforme al texto aplicable en 2026.}
\section{Conformidad y condición segura}
La aeronavegabilidad vincula conformidad con un diseño aprobado y condición para una operación segura dentro del régimen aplicable. La certificación inicial y el seguimiento durante la explotación son partes relacionadas de esa demostración. Una aeronave puede conservar un certificado y, aun así, tener una condición que impida volar hasta resolverla.

El certificado de aeronavegabilidad de una aeronave individual no sustituye al certificado de tipo del diseño. Su continuidad se relaciona con las condiciones de 21.A.181 y con las reglas de aeronavegabilidad continuada. No se concluye que un certificado emitido una vez permita operar indefinidamente sin requisitos posteriores.
\section{Qué es un ARC}
El certificado de revisión de la aeronavegabilidad (ARC) registra el resultado del proceso de revisión correspondiente. No es la liberación tras una tarea de mantenimiento ni una nueva certificación de tipo. La revisión incorpora documentación y examen físico conforme a las disposiciones aplicables; los registros y la condición real deben ser consistentes.

En la consolidación de 1321/2014 de agosto de 2026, M.A.901 establece el proceso general y M.A.903 desarrolla la revisión. M.A.901(b) distingue la emisión de un ARC y una recomendación a la autoridad según las condiciones que especifica; no se reproduce sin cambios el esquema de diapositivas de versiones anteriores.
\section{Qué debe comprobar la revisión}
La revisión documental debe permitir reconstruir matrícula e identificación, configuración, programa, mantenimiento, directivas, modificaciones, reparaciones y componentes controlados, entre otros elementos exigidos por el apartado aplicable. La revisión física contrasta la aeronave y su condición con los datos y registros.

No se resuelve una discrepancia física diciendo que «en la hoja está correcto». Tampoco se acepta un dato que no puede atribuirse a la aeronave o componente correspondiente. El objetivo no es crear un segundo historial paralelo que sustituya al original, sino revisar la evidencia y gestionar discrepancias según procedimiento.
\ilustra{revision-arc}{Relación entre revisión documental, examen físico, discrepancias y resultado de revisión. Esquema de estudio sujeto al ámbito y condiciones de M.A.901 y M.A.903.}
\section{Emisión, validez y condiciones}
M.A.902 regula validez del ARC; debe leerse junto con M.A.901 y con las disposiciones de la organización y autoridad. Una extensión de validez no se considera automática por aparecer una casilla disponible en el formulario. Se comprueban las condiciones y privilegios concretos.

El texto de M.A.901(f) permite adelantar la revisión 90 días o menos antes de caducidad sin perder continuidad del patrón, en las condiciones del apartado. No significa que cualquier tarea de mantenimiento o cualquier prórroga tenga una tolerancia general de 90 días. Esa distinción evita convertir una flexibilidad procedimental específica en un margen universal de operación.
\section{AD, modificaciones y reparaciones}
Las AD aplicables deben cumplirse según M.A.303 o ML.A.303, en su caso. Los datos de modificaciones y reparaciones se tratan, entre otros puntos, en M.A.304. La revisión necesita conocer su estado: una modificación puede cambiar aplicabilidad de una AD o introducir instrucciones adicionales.

Un boletín cumplido puede formar parte de la evidencia de cumplimiento de una AD, pero se comprueba la revisión y la acción que la directiva exige. La anotación «SB realizado» no demuestra por sí sola haber satisfecho una inspección repetitiva o una condición posterior de la AD.
\section{Registros y transferencia}
M.A.305 y las disposiciones correspondientes a otros regímenes organizan el sistema de registros de aeronavegabilidad continuada. La trazabilidad de horas, ciclos y vida de componentes evita atribuir al estado individual una vida que no se ha demostrado. La transferencia de registros debe conservar continuidad y legibilidad.

Una falta de documentación no prueba automáticamente que una pieza esté físicamente dañada, pero puede impedir demostrar su condición de cumplimiento. Reconstruir un historial requiere evidencia aceptable y no se sustituye por asignar cero a un contador desconocido.
\caso{M16. Certificados que responden a preguntas distintas}{Después de una intervención se emite un certificado de puesta en servicio para el mantenimiento ejecutado. El propietario pregunta si eso renueva automáticamente el ARC. La respuesta académica es no: el CRS identifica la certificación de la intervención; el ARC corresponde a un proceso de revisión con requisitos y privilegios específicos. Puede haber relación documental entre ambos, pero no equivalencia.}
\section{Cómo resolver un caso de examen}
Se identifica régimen M o ML, estado del certificado, fechas, quién gestiona, organización y privilegios, historial, mantenimiento y discrepancias. Se localizan los puntos pertinentes de la versión aplicable y se formula una conclusión condicionada a la evidencia. Si falta la configuración o un dato documental, esa falta se registra expresamente, sin completar la historia con supuestos ocultos.
\section{ARC: revisión, expedición y extensión son actos distintos}
La revisión de aeronavegabilidad combina documentación y examen físico; la emisión del ARC requiere satisfacción de condiciones y tratamiento de hallazgos; la extensión de su validez es un procedimiento condicionado. No equivalen al CRS tras mantenimiento ni al certificado de tipo.

Las diapositivas docentes sobre entorno controlado y emisión por recomendación son un esquema histórico y se contrastan con los artículos vigentes. En Parte ML existen reglas específicas: en supuestos previstos, el adelanto de la extensión del ARC puede alcanzar 30 días sin perder el patrón, mientras el adelanto de la revisión física puede alcanzar 90 días. No se transfieren esos plazos a inspecciones AMP ni se presumen iguales en Parte M. Un ARC anterior no acredita por sí mismo que una AD pendiente haya sido satisfecha.


\section{Relacionar documentación y estado físico}
Una revisión de aeronavegabilidad no se reduce a fechas. Identidad, configuración, estado físico y cumplimiento deben reconciliarse. Un historial aparentemente completo puede omitir una instalación; un buen aspecto exterior puede coexistir con una limitación vencida. Una inspección general tampoco sustituye todos los registros.

M.A.901 y M.A.903 de la consolidación de 2026 permiten distinguir condiciones de emisión o recomendación y proceso de revisión; M.A.902 trata validez. Los privilegios y documentos dependen del escenario. No se atribuye la misma facultad a cualquier organización por trabajar habitualmente con la aeronave. ML necesita su propio análisis.
\section{Documentos con preguntas distintas}
Certificado de tipo, CRS y ARC tienen objetos diferentes. La existencia de uno no renueva automáticamente otro. Un CRS tras una tarea no añade por sí solo otro periodo de ARC; un ARC no demuestra correcta ejecución de cualquier mantenimiento posterior. Certificado de aeronavegabilidad y ARC también se distinguen.

Una fecha en planificación no constituye una decisión de autoridad. Una ventana permitida para anticipar revisión se usa en su escenario exacto; no se convierte en tolerancia general de mantenimiento o límites de vida.
\section{Resolver una discrepancia de estado}
Si una hoja dice «cumplido» sin evidencia, se identifica tarea, fuente, revisión, configuración y registro que sostendría el estado. Una inspección visual genérica no equivale a una acción específica de AD. Se registra el hallazgo y se aclara por el proceso pertinente, sin fabricar un cierre con datos supuestos.

«Pendiente de demostrar» puede ser la conclusión correcta aunque la fecha esté bien calculada. Localiza qué prueba falta, sin convertir ausencia de evidencia en prueba automática de una condición física concreta.

\Needspace{.75\textheight}
\section{Vuelos de comprobación y autorizaciones de vuelo}
Un vuelo posterior al mantenimiento puede servir para comprobar funciones que no se evalúan completamente en tierra. La necesidad se define por los datos y el proceso; M.A.301(i) incluye vuelos de comprobación de mantenimiento cuando sean necesarios. Esto no significa que todo mantenimiento obligue a un vuelo, ni que cualquier vuelo de prueba pueda realizarse sin analizar condición de aeronave, operación y tripulación.

Parte 21, subparte P, regula permisos de vuelo. El supuesto de 21.A.701(a) se refiere a aeronaves que no cumplen o no han demostrado cumplir requisitos aplicables, pero son capaces de vuelo seguro bajo condiciones definidas y para los fines enumerados. Entre esos fines figura trasladarse donde se realizará mantenimiento o revisión; también hay un supuesto para localizar averías o comprobar sistemas tras mantenimiento. El propósito por sí solo no demuestra que todos los requisitos de emisión estén satisfechos \cite{EU748}.

El punto 21.A.708 trata condiciones de vuelo, incluida configuración, restricciones, procedimientos, programa de ensayos cuando proceda y disposiciones de aeronavegabilidad continuada. 21.A.710 regula su aprobación y 21.A.711 la emisión del permiso según autoridad o privilegios de organización pertinentes. Aprobar condiciones y emitir permiso son pasos relacionados con responsabilidades propias. No se atribuye esa facultad a cualquier taller por tener Parte 145.

Un permiso no transforma una aeronave en plenamente conforme para toda operación. Conserva finalidad, condiciones y restricciones; cambios de configuración pueden afectar su validez y tratamiento. En un estudio se identifica primero por qué se necesita el vuelo, qué condición tiene la aeronave y qué régimen y datos se aplican. Después se determina la ruta competente, sin improvisar un perfil de ensayo ni usar el libro como autorización.
\caso{M16-B. Traslado académico a un taller}{Una aeronave ficticia tiene una discrepancia y se propone volar al lugar de mantenimiento. Antes de calcular combustible o fecha se pregunta si está en condición de volar bajo el régimen ordinario o si necesita otra vía. La mera conveniencia del traslado no acredita seguridad ni concede permiso. Se reúnen configuración, discrepancia, datos y finalidad para analizar condiciones y aprobación en su marco.}

\repaso{¿Puede una ventana de revisión ampliar una inspección del programa?}{No se deduce tal autorización. Cada disposición tiene objeto y condiciones; se verifica el requisito de la tarea.}

\resumen{Tipo, certificado individual, ARC y CRS son documentos diferentes. La revisión compara documentos y aeronave. Una flexibilidad de un punto no es una tolerancia general. La falta de trazabilidad no se corrige reiniciando contadores ni declarando cumplimiento sin evidencia.}
\fuentes{Reglamento 748/2012, 21.A.181 \cite{EU748}; 1321/2014, M.A.303--305 y M.A.901--903 \cite{EU1321}; modificación 2026/100 \cite{EU2026-100}; material docente de aeronavegabilidad y ARC.}
